\documentclass[../main.tex]{subfiles}
\begin{document}
%==========================================TIÊU ĐỀ TẬP========================================
\begin{table}[h]
    \begin{adjustwidth}{-1cm}{}
        \begin{tabular}{>{\centering\arraybackslash}p{6cm}|>{\raggedright\arraybackslash}p{12cm}}
            \multirow{5}{6cm}
            % Cột bên trái: ÉPISODE và số tập
            {\\[-40pt]
                \humanist\fontsize{20pt}{20pt}\selectfont \centering
                \color{eptype}FOLGE\color{black}\\
                \houseterror\fontsize{72pt}{72pt}\selectfont
                \color{epnum}4\color{black}\\[6pt]
                \cafeteria\fontsize{20pt}{20pt}\selectfont\color{epnum}(M1-4)\color{black}
                \\[18pt]
                \humanist\fontsize{18pt}{18pt}\selectfont
                \color{horror} GRUSEL?\\
                \humanist\fontsize{14pt}{14pt}\selectfont
                \color{epattr}WILLKOMMEN, \colorbox{vol}{\color{Knight} Temma}!
                \color{black}
            }
             &
            % Dòng thứ nhất: tên tập tiếng Nhật
            {\fotskip\fontsize{18pt}{18pt}\selectfont \ruby{大人}{おとな}がビビってるところ\ruby{見}{み}て}                                                                                                                                                                         \\[6pt]
            % Dòng thứ hai: tên tập tiếng Anh
             & {\cafeteria\fontsize{18pt}{18pt}\selectfont Courage}                                                                                                                                                                               \\[4pt]
            % Dòng thứ ba: tên tập tiếng Việt
             & {\vnmsans\fontsize{18pt}{18pt}\selectfont Lòng can đảm}                                                                                                                                                                         \\[8pt]
            % Dòng thứ tư: Nhân vật xuất hiện
             & {\caviar{\fontsize{14pt}{14pt}\selectfont Nhân vật: \emp{kaigou2} \emp{naomi2} \emp{game} \emp{astel} \emp{temma}}}\\
            % Dòng thứ năm (nếu có): Nhân vật xuất hiện trong Omake
            & {\abv Truyện phụ:} \emp{kuon2} 
            \\[8pt]
            % Dòng cuối cùng: Thời gian diễn ra của tập (thường sẽ là ngày phát hành)
             & {\continuum\fontsize{16pt}{16pt}\selectfont Ngày phát hành: 18/02/2022}                                     \\[18pt]
        \end{tabular}\\[8pt]
    \end{adjustwidth}
\end{table}
%=========================================TRUYỆN CHÍNH========================================
\color{black}

\txtbx[2]{
    {\color{vol} \uline{Tóm tắt:}} Temma chơi game kinh dị đến mức hoảng loạn chạy vào văn phòng, giật mình vì tiếng bước chân của chính mình, tiếng nước sôi và lò vi sóng. Kaigou và Naomi phải dỗ dành mãi mới bình tĩnh được. Cuối cùng cả nhóm phát hiện Astel đang trốn trong tủ lạnh. Kuon về với mặt lem vữa xi măng, kể chuyện Mio từng tè ra quần vì bị dọa khiến cả phòng cười nghiêng ngả.
}

Trời dần sẩm tối, ánh hoàng hôn cuối cùng cũng tắt hẳn sau chân những tòa nhà cao tầng. Căn mancave của nhà \hls\ chìm trong bóng tối, chỉ còn lại ngọn đèn bàn lập lòe, thứ ánh sáng vàng ấm cố gắng chiếu xuyên qua đống sổ sách thu chi dày cộp xếp thành núi trên bàn làm việc.

Naomi ngồi co ro trên ghế, đôi mắt thỉnh thoảng lại lướt qua mớ giấy tờ mà Kaigou phải đắm chìm từ chiều. Lý do cô bé tám tuổi có mặt ở đây giờ này, tất cả cũng vì anh trai mình sau khi kết thúc ca làm tại văn phòng \hll, phải tất tả chạy sang công trường khu Ký túc xá. Cậu phải trộn hồ, xây dựng phòng ở mới cùng Ryujin và Hijaki, không muốn em gái phải hít bụi đá nên đã gởi gắm Naomi cho Kaigou trông nom vài tiếng. Suzuno và Ryuuhou thì vẫn còn đang ở trường, chưa về đến nhà.

Naomi đang cầm cây bút màu, cố gắng hoàn thành bức tranh tô dở trên tờ giấy trải ra ghế, thì bụng bỗng kêu lên một tiếng *ỌT ỌT* nho nhỏ, vang lên giữa sự im lặng của căn phòng. Cô bé lập tức rụi mắt, hai tay vắt vẻo kéo nhẹ vạt áo Kaigou, giọng nói nhỏ như tiếng muỗi kêu: ``Kaigou-oneechan\dots\ bụng em kêu rồi này. Khi nào bọn mình mới được đi ăn tối vậy ạ?''

Kaigou vẫn đang đập trán liên tục xuống đống tài liệu, hai má phồng lên như hai quả bóng xà phòng. Nghe tiếng bé cô em út nói, cô chỉ có thể thở dài mệt mỏi, toàn thân như sắp tan chảy ra ghế: ``Mừ'' Tại sao lại nhiều việc thế này không biết\dots''

Ngay lập tức, từ chiếc đồng hồ thông minh đặt giữa bàn, giọng nói lạnh lùng của Game vang lên, như một cơn gió lạnh xuyên qua căn phòng ấm áp: ``Nếu cô không lướt mạng xã hội cả chiều nay, thì giờ này đã xong hết và đưa Naomi-user đi ăn tối rồi. Cớ gì lại phải ngồi than thở ở đây?''

Kaigou lập tức nhảy dựng lên, hai tay vắt vào hông như chuẩn bị cãi lộn: ``Tôi có nghĩ đâu rằng mình phải làm thêm đến tận giờ này chứ! Mà Daidou-san đâu rôi? Sao lại vứt nguyên đống sổ sách này cho tôi một mình mà không giúp đỡ gì thế!''

Game vẫn bình thản, giọng nói chẳng có chút xúc động nào: ``Giờ cô nhắc tôi mới nhớ\dots\ Ông ấy đã biến mất không dấu vết từ trước khi tôi bắt đầu làm việc cùng cô rồi.''

Kaigou bĩu môi, vẻ mặt bực bội hiện rõ lên khuôn mặt. Cô lẩm bẩm một mình, nhưng đủ to để Game nghe thấy: ``Gã ấy thậm chí còn chẳng bao giờ xuất hiện từ đầu! Đàn ông gì mà chẳng ga lăng tí nào\dots\ Chả bù cho Aniki của tôi đâu.''

Game không chịu thua, liền chêm thêm một câu như muốn chọc tức Kaigou: ``Có cần tôi gọi điện cho Kuon-user bảo cậu ấy ngưng việc ở công trường mà chạy về đây phụ cô không?''

Nghe vậy, lòng tự trọng của Kaigou lập tức trỗi dậy. Cô vỗ ngực dõng dạc, như một nữ tướng chuẩn bị ra trận: ``Khỏi cần! Một mình tôi cũng giải quyết xong hết đống này! Tôi đâu có yếu đuối đến mức phải phiền phức nhờ vả người khác!''

Game chỉ có thể thở dài một tiếng, tiếng thở dài vang lên từ loa đồng hồ: ``Giá như cô cũng nói câu đó từ đầu giờ chiều. Thì bây giờ đâu phải tự làm khó mình, lại còn để Naomi-user phải ngồi nhịn đói theo cô nữa.''

Kaigou lập tức cau mày, chỉ tay thẳng vào chiếc đồng hồ, giọng nói đầy bực bội: ``Này! Tôi không có điếc đâu nha! Ông ngậm cái loa lại giùm đi!''

Naomi thấy vậy, liền xoa xoa cái bụng đang xẹp lép, hai mắt long lanh như ngọc nhìn Kaigou, giọng nói nịnh nọt: ``Kaigou-oneechan\dots\ Nhưng mà em đói thật rồi mà\~{} Lúc chiều Onii-chan có dặn, nếu anh về muộn thì chị Kaigou-oneechan sẽ dẫn em đi ăn món gì ngon lắm đó\dots''

Kaigou lập tức mềm nhũn ra trước vẻ mặt đáng yêu của bé Naomi. Cô vội vàng xoa đầu cô bé, giọng nói nhẹ nhàng an ủi: ``Được thôi, được thôi. Để chị hoàn thành nốt mấy tờ này thôi, rồi chị dẫn em đi ăn ramen ở quán cuối phố nhé? Quán đó có thêm kem dâu mà em thích lắm.''

Trong lúc Kaigou và Game còn đang cãi nhau chí chóe, từ phía hành lang bên ngoài bất ngờ vang lên tiếng bước chân chạy dồn dập, gõ chan chát xuống sàn nhà.

``Ai lại đến văn phòng vào cái giờ này nhỉ?''' nàng Tổng tò mò ngước đầu lên.

``\textbf{Lần này mình sẽ bị ám đến chết mất!!}''

Một giọng nam hơi cao cất lên đầy thảng thốt. *RẦM!* Cánh cửa văn phòng bị đẩy văng ra, một chàng trai tóc vàng khoác trên mình bộ trang phục hiệp sĩ xông thẳng vào trong, thở hồng hộc như vừa thoát khỏi hang hùm.

Naomi giật mình giật thót, nhanh như sóc tụt xuống khỏi ghế, chui tọt vào sau lưng Kaigou, hai tay ôm chặt lấy vạt áo chị gái, tròn xoe mắt thò đầu ra nhìn.

Kaigou trố mắt kinh ngạc: ``Có chuyện gì thế em, Temma-kun? Làm gì mà thở như đứt hơi vậy?''

\chrint

Người vừa đạp cửa lao vào văn phòng trong trạng thái hoảng loạn sắp ngất tới nơi là Kishido Temma (\ruby{岸}{き}\ruby{堂}{しど}\ruby{天}{てん}\ruby{真}{ま}, \ruby{Ki-shi}{Ngạn}-\ruby{đô}{Đường}\ \ruby{Ten}{Thiên}-\ruby{ma}{Chân}), thành viên thuộc Thế hệ thứ Hai của \hls. Với tư cách là một đệ tử tốt nghiệp từ Học viện Hiệp sĩ, cậu luôn mang đến năng lượng tích cực và nụ cười ấm áp với mọi người.

\model{25}{r}{4.5cm}{temma}

Từ nhỏ, con đường nghĩa hiệp đã in sâu vào máu của Temma, và ngay sau khi tốt nghiệp, cậu đã bước chân ra thế giới rộng lớn để tự mình trải nghiệm cuộc sống. Lần đầu tiên nếm thử món ramen đường phố, hương vị đậm đà đã khiến cậu kinh ngạc đến mức không thể nói nên lời. ``Thế là món ăn phổ biến nhất Nhật Bản lại có sức hấp dẫn khó cưỡng đến thế?'' Cậu đã thốt lên như vậy sau lần đầu thưởng thức, làm Kaigou đứng cạnh phải bật cười thành tiếng.

Có thể nói, tính cách của Temma chính là một sự tương phản kỳ lạ, hoàn toàn không đi liền với hình ảnh bên ngoài. Khi bước vào thế giới của các trò chơi sinh tồn, chàng hiệp sĩ luôn mỉm cười kia lập tức biến thành một kẻ chơi game đầy ắp chiến lược sắc bén, thậm chí có phần ``ác nghiệt'' trong trò chơi''. Mỗi khi thành công ``giết'' được ai đó, cậu luôn cười vô tư và chẳng hề hối tiếc. Chính điều này đã khiến người hâm mộ và các thành viên \hls\ khác đặt cho cậu biệt danh ``Hiệp sĩ Bóng tối''.

Bên cạnh đó, khi không ở trong trạng thái chơi game hỗn loạn, Temma lại là một người rất thực tế và luôn nhìn nhận vấn đề một cách tích cực. Cậu từng nhắn nhủ với người hâm mộ rằng dù rất trân trọng sự ủng hộ của mọi người, điều quan trọng nhất là hãy luôn chăm sóc và yêu thương bản thân mình trước tiên.

Trong một buổi stream chia sẻ về cuộc sống thường nhật, cậu cũng từng an ủi người xem rằng việc duy trì thói quen tập thể dục mỗi ngày có thể mang lại nhiều niềm vui bất ngờ. Một lần khác, Temma thẳng thắn chia sẻ rằng cậu không bao giờ quá bận tâm đến số lượt xem của mỗi buổi phát trực tiếp. ``Khi mình stream, mình chỉ chia sẻ những gì mình thích thôi, chấp nhận rằng người xem có thể sẽ thích hoặc không thích nội dung đó,'' cậu đã nói như vậy với một nụ cười chân thành.

Bên cạnh đó, Temma còn tạo dựng cho mình hình ảnh một ``thiên tài'' (tensai) nhờ những cách thức độc đáo để làm sôi động các buổi phát trực tiếp của mình, từ những tình tiết ngẫu nhiên, kỳ lạ cho đến những cú twist hài hước không ai ngờ tới. Chính vì vậy, cả các thành viên HOLOSTARS lẫn người hâm mộ đều đồng ý rằng đây chính là điểm làm nên sự độc đáo của Temma—cách cậu nhìn nhận mọi chuyện thường khác biệt với thông thường, và đôi khi khiến người khác khó hiểu. Ở văn phòng, mọi người luôn trông chờ vào những ý tưởng bất ngờ mà cậu mang lại.

Về ngoại hình, cậu là một chàng công tử với mái tóc vàng ngắn gọn, đôi mắt xanh lam trong trẻo như bầu trời ngày nắng. Dù trong mô hình 3D có vẻ hơi mảnh khảnh, nhưng sức mạnh thực tế của Temma lại không phải chuyện đùa: cậu sở hữu đai đen karate, đủ mạnh để bế kiểu công chúa bất kỳ bạn cùng thế hệ hay cõng cả Kageyama Shien trên vai đi đâu cũng thoải mái.

Trang phục thường ngày của cậu cũng thể hiện rõ nét trang trọng của một hiệp sĩ: chiếc áo khoác trắng cổ điển với hai hàng cúc hình thoi viền vàng, mặc bên trong áo sơ mi đen lịch sự. Cậu phối cùng quần tây trắng, cà vạt xanh có kẹp cà vạt, dây đeo đùi và đôi giày da đen bóng lộn, tạo nên một hình ảnh vừa nghiêm túc vừa không kém phần thanh lịch.

\chrint

Nhưng chàng hiệp sĩ nào có còn tâm trí để ý tới lời hỏi thăm của Kaigou nữa. Cậu đứng huơ huơ giữa phòng, tay chân cứ cuống quýt như vừa trốn thoát khỏi đàn sói, giọng run rẩy tới mức chữ chữ cũng rớt ra: ``Mọi người ơi, tôi thề là cái game này--''

\img{m1-4a}

Lời còn chưa kịp hết miệng, Temma bỗng khựng lại. Cậu chớp chớp mắt, nhìn quanh căn phòng tối thui chỉ thấy mỗi Kaigou tóc đuôi ngựa và một bé gái nhỏ đang nhìn chằm chằm vào mình.

Cậu chàng Hiệp sĩ ngẩn người ra: ``Ể? Ở đây chỉ có mỗi chị thôi sao, Kaigou-senpai?''

Kaigou thở dài một hơi mệt mỏi, giọng vừa buồn cười vừa lo lắng: ``Mọi người về hết từ mấy tiếng trước rồi. Cậu có chuyện gì mà chạy tới đây hô hoán như bị ma đuổi thế?''

Naomi từ sau lưng Kaigou thò đầu ra, đôi mắt to tròn long lanh, giọng trong trẻo như tiếng chuông gió hỏi dồn: ``Anh ơi, anh bị ma đuổi thật ạ? Ma có răng nhọn không anh?''

Temma giật mình cúi xuống nhìn bé gái, lưng vẫn còn ươn ướt mồ hôi lạnh, vừa vỗ ngực vừa nói trong tiếng thở hổn hển: ``Thì\dots\ em vừa mới chơi một trò chơi kinh dị ở phòng bên cạnh mà! Cái game đó nó—''

*RẦM!*

Chưa kịp dứt câu, cánh cửa gỗ đằng sau lưng cậu bỗng dưng đóng sầm lại với một tiếng động lớn mà không có ai chạm vào. Tiếng vang dội cứ thế lơ lửng giữa căn phòng im lặng như tờ.

``\textbf{Eek!! C-Cánh cửa\dots!!}'' Temma giật bắn người, nhảy nhỏm lên đồng thời tay phải vội vã lôi thanh kiếm được giắt từ hông ra, chĩa thẳng về phía cánh cửa vừa đóng sầm, chân tay còn lại rụt lại như vừa thấy ma hiện ra.

\gif{m1-4b}{1}{69}

Naomi cũng giật mình ôm chặt eo Kaigou, nhưng rồi cô bé chớp chớp mắt, chú ý tới chiếc đồng hồ thông minh trên bàn đang nhấp nháy ánh sáng xanh dịu nhẹ.

Game từ màn hình đồng hồ cất giọng thản nhiên, như chỉ vừa mới làm một việc vô cùng vặt vãnh: ``À, xin lỗi nhé. Tôi vừa thử cơ chế cửa tự động mới cài đặt thôi. Mà này, hiệp sĩ kiểu gì mà nhát như thỏ đế thế?''

Temma mặt bừng đỏ lên vì ngượng, cố gắng vớt vát chút sĩ diện còn sót lại của một hiệp sĩ: ``K-Không phải! Chỉ là tôi giật mình vì âm thanh bất ngờ thôi!''

Game chẳng mấy tin tưởng, giọng vẫn hờ hững như đang xem một vở hài kịch: ``Yeah, yeah. Tùy cậu nói thôi. Cậu với Shirogane-san bên \hll\ đúng là một cặp bài trùng nhút nhát y hệt nhau.''

Kaigou vội vỗ vai Temma để ngăn cậu tiếp tục bị Game trêu chọc: ``Game này, đừng có phũ với người ta như vậy chứ\dots\ Temma-kun, để chị đi bật đèn lên cho đỡ tối nhé.''

*CẠCH!*

Công tắc đèn được bật lên, ánh sáng trắng của đèn trần lập tức xua tan hết bóng tối ảm đạm của căn phòng. Temma thở phào một hơi dài, như thể vừa thả được tảng đá nghìn cân ra khỏi ngực. Để bình tĩnh lại, cậu đi tới khu bếp nhỏ, rót nước vào bình siêu tốc rồi bật công tắc đun. Xong xuôi, cậu lại lững thững đi về ghế, chống cằm lên hai bàn tay, ánh mắt trống rỗng nhìn vào khoảng không.

\img{m1-4c}

Kaigou chống hai tay lên hông, nhìn cậu chàng với vẻ nghi ngờ: ``\dots Chị hỏi thật đấy, rốt cuộc là có chuyện gì khiến em hoảng hồn vậy?''

Temma thở dài một tiếng nặng nề, như thể vẫn chưa thoát khỏi không khí kinh dị của trò chơi vừa rồi: ``Em vừa chơi một trò kinh dị. Nó thế nào nhỉ\dots''

Game lập tức chêm vào, giọng đầy vẻ soi mói như đã đoán ra từ trước: ``Kiểu mấy trò dùng pha hù dọa rẻ tiền (jumpscare) ấy á? Nếu là thế thì tôi thấy nó nhạt nhẽo cực kỳ, chẳng có gì đáng sợ cả.''

Chàng Hiệp sĩ lắc đầu quầy quậy, phủ nhận lời của Game: ``Không phải cái trò rẻ tiền đó! \dots À mà, chắc cũng có một chút, nhưng mà nó khác lắm! Nó\dots\ như thế nào ấy nhỉ\dots''

Kaigou tặc lưỡi một tiếng, đoán thử dựa trên những trò chơi kinh dị cô đã từng trải qua: ``Chắc lại là mấy cái trò giải đố tìm đồ vật trong căn nhà ma u ám chứ gì? Cái đó chị chơi suốt mà, chẳng có gì đáng sợ.''

Temma gãi đầu gãi tai, cố gắng diễn tả cảm giác khó tả mà trò chơi vừa để lại trong mình: ``Cũng chẳng phải\dots\ Nó cứ kỳ kỳ thế nào ấy\dots''

Cậu cứ ngồi đó, vẻ mặt chán nản và bất an lan tỏa khắp người. Cả phòng rơi vào im lặng, không ai nói thêm lời nào, duy chỉ có\dots

*SÙNG SỤC''*

Tiếng bình đun nước bắt đầu sôi ùng ục, hơi nước nóng bốc lên nghi ngút, che mờ một góc căn bếp. Đúng lúc nước sắp sôi sục lên đỉnh điểm, Temma bỗng giật phắt người ngồi dậy. Cậu lao như một mũi tên tới góc bếp, ấn cái rắc tắt phụt chiếc ấm siêu tốc, miệng lẩm bẩm trong sự sợ hãi vô hình vẫn chưa tan: ``Không ổn rồi\dots!''

Kaigou và Game đồng loạt đổ mồ hôi hột trên trán, không hiểu cậu lại bị gì nữa. Hai người đồng thanh lên tiếng, cùng một câu hỏi: ``Nước đã sôi hẳn đâu em/anh ơi\dots''

Temma cười gượng gãi gãi cổ, gương mặt vẫn còn tái nhợt vì hoảng sợ, cố gắng biện minh cho hành động kỳ lạ của mình: ``E-Em thấy nước như thế này cũng\dots\ đủ nóng rồi mà\dots''

Nói rồi, cậu vội vã xé gói bột ca-cao đổ vào cốc, rót dòng nước ấm từ trong bình vào rồi dùng thìa khuấy lấy khuấy để. Toàn bộ thao tác đều được Temma thực hiện với tốc độ ánh sáng, như thể có một thế lực tâm linh nào đó đang đuổi theo cậu để truy sát.

Khuấy xong ly ca-cao nửa nguội nửa nóng của mình, Temma ngước lên, giọng vẫn còn run rẩy nhẹ khi hỏi Kaigou: ``Chị muốn làm một cốc ca-cao cho ấm người không, Kaigou-senpai?''

Nàng Tổng gật đầu, vẫn còn nhìn cậu với vẻ không thể tin được: ``Ừ, làm cho chị một cốc đi.''

Naomi đứng bên cạnh kiễng chân, giơ bàn tay nhỏ xíu lên nịnh nọt, muốn có một ly cho riêng mình: ``Anh hiệp sĩ ơi! Làm cho em một cốc nữa với ạ! Em cũng muốn uống ca-cao nóng\dots\ Với lại, tay anh đang run kìa!''

Temma giật mình nhìn xuống cô bé 8 tuổi ngây thơ đang đứng trước mặt, mặt bỗng hơi đỏ lên vì ngượng khi bị phát hiện ra sự sợ hãi của mình: ``A, có cả bé Naomi-chan ở đây nữa sao? Đ-Đâu, anh đâu có run! Được rồi, anh làm cho hai chị em mỗi người một cốc nhé!''

Cậu lẩy bẩy lấy thêm một chiếc cốc sứ trên chạn xuống, hai bàn tay run rẩy đến mức gốm sứ va vào nhau tạo nên những tiếng *LENG KENG* nho nhỏ.

Kaigou nhìn cậu chàng bằng ánh mắt đầy lo ngại, không nén nổi thở dài: ``Em có thực sự ổn không đấy Temma-kun? Mặt em xanh như đít nhái rồi kìa.''

Temma giật bắn người, giấu hai tay ra sau lưng, gượng cười: ``K-Không có gì đâu chị! Em hoàn toàn bình thường!''

Game không nghĩ như thế, nhưng ông ấy chỉ giữ im lặng trên màn hình đồng hồ. ``\textit{Đến cả pha cốc ca-cao mà Kishido-chan trông cũng hấp tấp và run như cầy sấy,}'' ông ta lẩm bẩm quan sát.

Cậu chàng Hiệp sĩ đưa cốc ca-cao nguội ngắt lên miệng uống một ngụm lớn, rồi ngay lập tức nhăn nhó mặt mày: ``E-Em ổn! *ừng ực\dots* Oái! Lạnh quá!''

Naomi đứng bên cạnh, thản nhiên thổi phù phù vào cốc ca-cao nóng hổi của mình, ngước mắt nhìn Temma như nhìn một sinh vật lạ: ``Thì lúc nãy anh tắt bình siêu tốc trước khi nước sôi mà. ''

Game chệm thêm vào: ``Bé Naomi nói đúng đấy. Nước đã sôi hẳn đâu mà đòi nóng.''

Tức tốc, chàng Hiệp sĩ vội vã bê cốc ca-cao nguội tống vào lò vi sóng gần đó rồi ấn nút quay: ``V-Vậy tôi sẽ quay cho nó nóng lên!!''

\img{m1-4d}

Thấy bộ dạng hoảng loạn mất kiểm soát của cậu, Kaigou và Game đều dè chừng lắc đầu ngán ngẩm: ``\textit{Thật sự cậu ta có ổn không đây\dots}''

Không chịu ngồi yên, Temma bắt đầu đi vòng quanh khắp căn phòng làm việc tĩnh lặng, dáng vẻ bồn chồn lo âu cực độ: ``Tại sao mọi người lại về sớm như vậy chứ\dots\ Để lại một mình em ở đây\dots''

*CỘP CỘP* *CỘP CỘP*

Từng tiếng bước chân vang lên đều đặn trong không gian vắng lặng. Bất chợt, cậu chàng khẽ giật mình gẩy người một cái, vô thức ngoái đầu lại phía sau gào lên: ``\textbf{Tiếng gì vậy!? Có phải là chị không, Kaigou-senpai!?}''

Kaigou lắc đầu: ``Chị im lặng ngồi đây nãy giờ mà.''

Temma lại bước tiếp vài bước. *CỘP CỘP* *CỘP CỘP*

Chàng Hiệp sĩ lại giật thót người nhảy chổm lên, mắt trợn ngược ngoái lại: ``\textbf{Ai vậy!? Ai vừa đi đằng sau tôi!? Có phải là ông không, Game-san!?}''

\img{m1-4e}

Naomi tay cầm cốc ca-cao, đứng kiễng chân chỉ thẳng vào đôi đế giày giáp kim loại của cậu, hồn nhiên cất giọng dập tắt cơn hoảng loạn: ``Temma-oniichan! Đấy là tiếng giày của anh dẫm xuống sàn nhà mà! Mỗi lần anh bước đi là nó kêu cộp cộp đó, đâu có ai đi đằng sau anh đâu!''

Game cũng lắc đầu: ``Tôi không xấu tính đến mức rảnh rỗi đi dọa cậu đâu.''

Dù được bé Naomi chỉ rõ nguyên do, Temma vẫn không thoát khỏi ``cơn mộng mị''. Cậu tiếp tục cuộc du ngoạn hỗn loạn khắp phòng. Càng hoảng, bước chân cậu càng dồn dập hơn.

*CỘP CỘP* *CỘP CỘP* *CỘP CỘP* *CỘP CỘP*

``T-Tiếng\dots\ Tiếng động này lại xuất hiện nữa rồi!!''

Cứ mỗi lần tiếng bước chân của chính mình vọng lại, Temma lại dừng khựng, giật mình nhìn đằng sau, rồi lại nhìn về phía Kaigou và Naomi. Cô em út lại chớp mắt nhắc lại không hề do dự: ``Lại là chân anh đó!'' và vòng lặp quái đản cứ thế tiếp diễn.

*CỘP CỘP* *CỘP CỘP* *CỘP CỘP* *CỘP CỘP*

Càng lúc, cậu chàng Hiệp sĩ hoảng loạn tột độ, chân cậu càng bước nhanh như chạy. Kaigou và Game bắt đầu đổ mồ hôi hột trước độ nhát gan tới độ khó tin của cậu

``\textbf{Ai đó!?}''' Temma gào lên.

Game gắt gỏng thẳng mặt: ``\textbf{Tiếng bước chân của cậu từ nãy tới giờ đấy đồ ngốc!!}''

*TINH!*

Đúng lúc này, chiếc lò vi sóng phát ra tiếng chuông báo hoàn thành. Ngay lập tức, Temma giật nảy mình, ôm đầu ôm mặt thét lớn: ``\textbf{Khôôôôôôông!!}''

\img{m1-4f}

Naomi giật mình mất một giây, rồi cô bé phì cười nhe răng khẩy: ``Lò vi sóng hâm xong ca-cao rồi kìa anh ơi! Kêu *tinh* một cái thôi mà anh cũng hét lên như bị ma cắn vậy!''

Kaigou vội vã xua tay trấn an: ``B-Bình tĩnh lại nào Temma-kun! Tiếng lò vi sóng thôi mà!''

Chàng tóc vàng run rẩy mở cửa lò, lấy cốc ca-cao ra thở hổn hển: ``T-Tưởng gì\dots\ Tiếng bước chân với lò vi sóng thôi sao\dots''

Cậu đi về phía bồn rửa mặt, cố gắng xốc lại tinh thần: ``Bình tĩnh lại nào Temma! Mình sẽ rửa mặt bằng nước lạnh một chút cho tỉnh táo\dots''

*XÈÈÈOOO!*

Vừa mở vòi nước, âm thanh dòng nước chảy mạnh đập vào lòng bồn inox vang lên. Temma lại phát hoảng nhảy dựng ngược lên trời, gào to: ``\textbf{Chỉ là tiếng nước thôi!!}''

\img{m1-4g}

Kaigou, Game và Naomi đứng trố mắt nhìn cảnh tượng hãi hùng trước mặt.

Cô em út nhà Hikawa thì̀ thầm vào tai chị mình: ``\hl{Kaigou-oneechan, anh hiệp sĩ này nhát hơn cả mấy bạn lớp ba ở trường em nữa\dots}''

Nàng Tổng nhìn cảnh tượng trước mắt chỉ biết thở dài ngán ngẩm: ``Em có vẻ dễ bị kích động quá mức rồi đấy Temma-kun\dots''

Game chêm vào phũ phàng: ``Tôi thấy cậu còn nhát cáy hơn cả những gì tôi có thể hình dung ra đấy. Chẳng hề xứng đáng với cái danh `Hiệp sĩ' mà cậu đang có cả.''

Kaigou tiến lại gần: ``Thôi bình tĩnh lại đi em! Để chị dìu em ra ghế ngồi.''

Cô nàng dìu chàng Hiệp sĩ đang hồn xiêu phách lạc xuống chiếc ghế sô-pha gần đó, bảo cậu tịnh tâm lại và hít thở sâu.

``Đáng sợ quá chị ơi\dots'' Temma nhắm tịt mắt, nhoài người dựa ngửa ra lưng ghế.

Game tặc lưỡi: ``Đúng là thảm hại. Còn thua cả Shirogane-PON bên \hll\ nữa kia.''

Kaigou nhăn mặt bênh vực: ``Đừng có so sánh cậu ta với Noel-tan như thế. Họ khác nhau hoàn toàn đấy.''

Temma bất ngờ lại giật thót người một cái nữa. Vị Trợ lý ảo tỏ vẻ khó chịu: ``Lại cái gì nữa đây!?''

Chàng tóc vàng đổ mồ hôi hột, hai mắt trợn tròn nhìn về phía chiếc tủ lạnh góc phòng đang mở hé một khoảng nhỏ. Cậu run rẩy chỉ tay: ``M-Mình có mở cái tủ lạnh này đâu nhỉ\dots?''

\img{m1-4h}

Game nhớ lại: ``Tôi không nhớ là cậu mở nó đâu. Chắc là Kaigou-user mở ra lấy đồ rồi quên đóng đấy.''

Kaigou lắc đầu quả quyết: ``Vô lý! Tôi nhớ là lúc nãy đã đóng nó chặt lắm rồi mà.''

Temma mồ hôi chảy ròng ròng: ``K-Kể cả thế, ta cũng không thể để nó mở hé đáng sợ như thế này được\dots''

``Đồng ý với Kaigou-user. Tốn tiền điện cả đấy chứ chẳng đùa đâu.''

Nàng Tổng đề xuất: ``\dots Hay sao ta không tiến lại kiểm tra xem bên trong có cái gì nhỉ?''

Cậu Hiệp sĩ nhìn cánh cửa tủ lạnh đang mở hé, người run cầm cập như cầy sấy: ``Nhưng mà em thấy sợ khi mở tủ kiểm tra lắm\dots\ Em không dám mở một mình đâu! Chị với Naomi-chan cùng làm với em nhé?''

Naomi hoàn toàn không biết sợ là gì. Cô bé hăm hở xắn tay áo, xung phong đi đầu: ``Được luôn! Để em mở cho! Trong tủ lạnh cùng lắm chỉ có bánh mì với hộp sữa thôi chứ làm gì có ma!''

Thấy một cô bé tám tuổi lại dũng cảm hơn cả mình, Temma ngượng chín mặt, cố gồng mình vặn lại: ``T-Tôi biết rồi! Tôi là hiệp sĩ, tôi không thể để một bé gái thấy mình yếu đuối được!''

Temma tiến lại gần tủ lạnh cùng hai chị em nàng Tổng. Cả ba đứng trước cánh cửa tủ hé mở.

``Một\dots''' Temma run run đếm.

``Hai\dots''' Kaigou tiếp lời.

``Ba!''' Naomi dùng hai tay kéo mạnh cánh cửa tủ ra.

*RẮC!*

Cánh cửa tủ lạnh mở tung. Tất cả đồng loạt khựng người, trố mắt nhìn thứ bên trong.

Ngồi lọt thỏm, co quắp trong ngăn mát tủ lạnh\dots\ lại chính là Astel. Chàng trai tóc xanh lơ ngơ ngác ngước nhìn ba người đứng bên ngoài.

\img{m1-4i}

Cậu chàng Bảnh bao xòe bàn tay rụt rè: ``Y-Yo\dots''

Kaigou ngơ ngác: ``Ê-tồ\dots\ Chào em?''

Game trên màn hình giật mình: ``Cái quái gì thế này\dots''

Naomi vừa nhìn thấy cậu bé trong tủ lạnh, hai mắt tức thì sáng rực. Em vỗ tay chộn rộn, ngón út nhỏ xíu chỉ thẳng vào trong tủ, giọng reo lên vui như vừa tìm được báu vật: ``A! Có một anh tóc xanh lơ đang ngồi biến thành kẹo kem ở trong tủ lạnh kìa!''

Lúc này, trái ngược với sự ngơ ngác của Kaigou và sự sốc của Game, Temma mới là người phản ứng mãnh liệt nhất. Cậu như kẻ sa mạc khát cỡi mấy tuần trông thấy nguồn nước, mặt mày sẵn niềm vui vô bờ bến lao tới.

Temma túm lấy cổ áo Astel, lôi cậu bạn ra khỏi ngăn mát tủ lạnh như kéo một món đồ quý giá rồi ôm chặt lấy không buông. Cậu reo lên trong giọng sung sướng như vừa thoát khỏi địa ngục: ``Người hùng của mình ơi! Astel ơi! Cậu cứu tớ rồi!''

Kaigou đứng bên cạnh vẫn chưa hết ngạc nhiên, vừa nhìn cặp bạn bè quấn quýt vừa nhìn vào cái tủ lạnh còn đang mở hé. Nàng Tổng ngơ ngác hỏi với giọng đầy thắc mắc: ``T-Tại sao em lại chui tọt vào trong tủ lạnh ngồi thế hả Astel?''

Astel vừa bị lôi ra khỏi không gian lạnh lẽo, cả người còn lạnh ngắt và đổ mồ hôi hột vì bí bách. Cậu gãi đầu cười gượng, cố gắng giải thích tình thế khó hiểu của mình: ``Em\dots\ Em chui vào đó trốn vì em nghe thấy cậu ta vừa chạy vừa hét vào văn phòng, sợ cậu ấy đuổi theo đòi ôm em\dots\ Ai ngờ thế nào em lại trở thành `người hùng' giải cứu cậu ấy luôn\dots''

\img{m1-4j}

Kaigou nghe xong câu giải thích chỉ biết đập tay mạnh lên trán. Nàng thở dài bất lực, giọng nói đầy sự bất lực trước tình huống vô lý này: ``Thiệt hả giời\dots''

Một bàn tay nhỏ lướt qua màn hình đồng hồ thông minh, Game thở dài xem vở hài kịch đang diễn ra trước mắt. Màn hình xanh nhấp nháy một câu chữ vui vẻ, như thể vừa ghi lại khoảnh khắc kỳ lạ nhất trong ngày: ``Nhật ký văn phòng hôm nay: Thu hoạch được một cá thể Astel Leda hoang dã trong tủ lạnh tại tầng 3!''

Sau vài giây im lặng, vị trợ lý AI liền buông một câu tổng kết cay đắng từ loa đồng hồ: ``Tôi chưa bao giờ thấy một hiệp sĩ nào lại dễ bị hú vía bởi chính tiếng bước chân của mình như thế này. Nhát cáy thua cả trẻ con!''

%==========================================TRUYỆN PHỤ=========================================
\omk

Bầu không khí trong văn phòng vẫn còn đang rung lên vì tiếng cười rúc rích của Naomi và những tiếng lắc đầu ngao ngán của Kaigou. Temma vẫn chưa hề buông người bạn thân của mình ra, ngược lại còn siết chặt thêm một chút, sợ rằng nếu buông tay thì cậu sẽ tan biến đi mất.

Cậu vẫn giữ nguyên tư thế ôm chầm lấy Astel, mặt cắm vào vai áo của người kia, giọng sụt sịt như đứa trẻ vừa bị trông ma: ``May quá\dots\ Có cậu ở đây rồi\dots\ Mình không còn sợ nữa!''' Temma sụt sịt, vẫn siết chặt lấy Astel như gấu ôm cây.

Cậu chàng Bảnh bao nhăn nhó, bất lực đẩy cái đầu vàng kè ra: ``Ờm\dots\ Temma ơi, cậu định ôm tôi đến bao giờ nữa vậy? Cổ tôi sắp gãy rồi đây này!''

Naomi đứng bên cạnh khoanh tay, che miệng cười rúc rích: ``Temma-oniichan ôm Astel-oniichan chặt như keo dán sắt luôn! Giống hệt em ôm con gấu bông lúc ngủ!''

Kaigou gãi đầu gãi tai khuyên can: ``A-nô\dots\ Temma-kun, em có thể thả em ấy ra được rồi đó.''

Temma lắc đầu quầy quậy: ``Không! Em muốn ôm cậu ấy thêm lúc nữa cho bớt sợ!''

Game trên màn hình phũ phàng cất giọng: ``Cậu biết là Leda-san đang ngộp thở tới mức sắp thăng thiên rồi không?''

Astel gật đầu lia lịa: ``Ông ấy nói chuẩn đấy! Thả tôi ra mau!''

``Rồi\dots\~{}''' Temma ngậm ngùi buông tay.

Nàng Tổng hắng giọng, quay sang hỏi: ``Mà sao em vẫn chưa về vậy Astel?''

Chàng trai tóc xanh lơ thở dài giải thích: ``Em cũng chuẩn bị về đây này. Nhưng ngay khi nghe thấy tiếng cậu ta vừa chạy vừa gào thét lao tới, em còn cách nào khác ngoài chui tọt vào tủ lạnh trốn đâu!''

Game cà khịa: ``Sao cậu không tìm chỗ nào khác mà trốn? Chui vào đó làm cậu Hiệp sĩ của chúng ta hú vía suýt rớt tim ra ngoài đấy.''

Astel gạt tay phân bua: ``Ông nhìn quanh xem còn chỗ nào trốn được nữa đâu! Tôi còn chẳng có ý định hù cậu ta nữa là! Mà Kaigou-senpai, sao giờ này chị vẫn còn ngồi đần mặt ra đây?''

Kaigou phồng má: ``Chị còn đống công việc thu chi chưa làm xong chứ sao\dots''

Game chệm một câu xanh rờn: ``Chính xác là vì cô ấy lướt mạng câu giờ từ đầu giờ chiều nên việc mới dồn lại một đống đấy.''

Nàng Tổng đỏ bừng mặt, đập bàn quát: ``\textbf{Tôi không có! Tại tôi không biết hôm nay sẽ phải làm thêm việc chứ bộ!}''

Trợ lý ảo tặc lưỡi: ``Cô làm việc còn thua cả Kuon-user đấy.''

Kaigou nghiến răng: ``Im đi cái ông này!''

Thấy màn cãi nhau chí chóe của nàng Tổng và vị trợ lý AI, Temma và Astel đều bật cười. Cậu chàng Hiệp sĩ tủm tỉm: ``Hai người đúng là ăn ý thật đấy.''

Kaigou và Game đồng thanh hét lên: ``\textbf{Bọn tôi không có!!}''

*CẠCH\dots*

Đúng lúc này, cánh cửa gỗ bất ngờ từ từ đẩy ra, kèm theo một tiếng gọi kéo dài thều thào: ``Hế-lô\dots''

Temma hoảng loạn tột độ, giật nảy người nhảy tót lên đu trên cổ Astel làm cậu bạn chới với. Kaigou cũng giật mình thót tim nhìn ra cửa.

Ở ngưỡng cửa, một bóng đen cao rỏng lù lù xuất hiện với gương mặt trắng bệch như ma. Temma gào lên chói tai, còn Kaigou run rẩy hỏi: ``A-Ai thế!?''

Naomi thì hoàn toàn không sợ. Cô bé chớp mắt nhìn kỹ rồi hăm hở chạy lại gần, tinh nghịch chỉ tay: ``A! Onii-chan kìa! Sao mặt anh trắng tinh như vừa nhúng đầu vào thùng bột mì thế kia?''

Bóng đen thở dài đánh thượt một cái, lấy tay quệt bớt lớp vữa khô trên mặt: ``Tôi đây\dots\ Hikawa Kuon đây\dots''

Thấy giọng nói quen thuộc của Kuon, Kaigou mới vuốt ngực thở phào nhẹ nhõm: ``Trời đất ơi! Hóa ra là anh à\dots\ Làm tốn mất mấy nhịp tim!''

Astel vẫy tay: ``A, Kuon-senpai!'' Kuon bước vào phòng, lắc đầu bực bội: ``Chào mấy đứa\dots\ *khẹc khẹc*\dots''

Kaigou tủm tỉm cười: ``Sao mặt anh trắng bệch ra thế kia, Aniki? Vừa mới đi spa dưỡng da về à?''

Kuon phũ phàng đáp: ``Dưỡng cái đầu em ấy. Anh vừa bên công trường trát gạch dọn dẹp phòng ký túc xá về, bị dính vữa xi măng đầy mặt còn chưa kịp rửa đây này!''

Nàng Tổng bật cười hề hề: ``Xin lỗi, xin lỗi nhé. Thế căn phòng mới xây đến đâu rồi?''

Cậu Tổng xoa đầu Naomi: ``Nhanh hơn anh tưởng. Thiếu mỗi sơn tường với lắp cửa nữa là hoàn thiện rồi. Còn em nữa Kaigou, sao giờ này còn chưa dắt Naomi đi ăn?''

Astel nhanh miệng đá đểu: ``Chị ấy ở lại vì làm thêm giờ đấy anh ạ. Lỡ nhịp từ chiều.''

Kuon gật gù: ``À\dots\ bị dồn việc hả?''

Nàng Tổng gãi đầu cười gượng: ``Thì cứ cho là vậy đi\dots''

Cậu Tổng thản nhiên đáp: ``Anh hiểu cảm giác của em mà. Làm nhiều mệt lắm.''

Astel bắt được sóng liền khoái chí trêu tới bến: ``Đúng là bạn trai của chị có khác, tâm lý và hiểu chị Kaigou ghê!''

Kaigou mỉm cười đỏ mặt, còn Kuon như thường lệ chỉ biết thở dài ngán ngẩm: ``Mày đừng có trêu anh nữa Astel\dots''' Rồi cậu quay sang vỗ vai Temma: ``Phải để cho mấy đứa thấy cảnh cơm chó vô lý này rồi, thông cảm nhé.''

Cậu chàng Bảnh bao hùa theo: ``Không sao, em thấy thích mà!''

Cậu Hiệp sĩ gãi cổ cười gượng: ``Mặc dù em cũng thấy hơi ghen tị với anh thật đấy Kuon-senpai\dots''

Kuon lắc đầu: ``Chẳng sung sướng gì đâu. Mà anh thấy Temma có vẻ nhát gan lắm đấy nhé.''

Temma giật thót người, mặt đỏ gay trong khi mọi người cười phá lên. Kuon cất giọng tiếp lời: ``Yên tâm đi, bên \hll\ cũng có một người thực sự tưởng mạnh mẽ nhưng hóa ra lại nhát cáy lắm.''

Temma tò mò ngước lên: ``Ai vậy anh?''

Kuon nhớ lại kỉ niệm cũ, bật cười phì: ``Nàng Sói đen Ookami Mio-chan chứ ai. Hồi anh mới vào làm, Shion-chan, Ryuuhou với anh có bày trò hù em ấy một trận nhớ đời\footnote{Xem lại {\flaregothic ÉPISODE 17}, quyển số Hai của \textit{holo no Graffiti}.}. Lần đó con bé còn nhập vai giả vờ làm nạn nhân ở lại với Haato-chan với Suba-chan nữa chứ. Mio-chan tưởng thật, sợ quá khóc thét rồi\dots\ tè cả ra quần luôn, đúng nghĩa đen đấy!''

Tất cả mọi người trong phòng lại được một trận cười nghiêng ngả. Game trên đồng hồ nhanh trẩu chêm vào: ``Thì ra cũng có người nhát như cậu đấy Kishido-san à!''

Temma nổi phừng phừng giậm chân: ``Đừng có xát muối vào lòng tôi nữa! Tôi biết là tôi nhát rồi mà!''

Naomi vỗ tay khoái chí: ``Hôm nay vui quá Onii-chan ơi!''

Một ngày thứ Sáu náo nhiệt cuối cùng cũng kết thúc. Astel và Temma lại rủ nhau sang phòng bên tiếp tục chinh phục trò chơi kinh dị, còn Kuon, Kaigou, Naomi và Game thì cùng nhau rảo bước trở về nhà sau một ngày làm việc vất vả và hỗn loạn một cách vô cùng bất đắc dĩ.

\end{document}